% --- Archon physics formalization source begin ---
% archon:physics
% archon:covers PhyXMiniProblems/problem_phyx_mini_0828.lean
% archon:source-report reports/phyx_mini/problem_phyx_mini_0828.source.json

\chapter{Physics problem phyx\_mini\_0828}
\label{ch:PhyXMiniProblems_problem_phyx_mini_0828}

\paragraph{Problem source.}
\#\# Physical scenario

Cleveland is \textbackslash{}(300\textbackslash{},\textbackslash{}mathrm\{miles\}\textbackslash{}) east of Chicago. A plane leaves Chicago flying due east at \textbackslash{}(500\textbackslash{},\textbackslash{}mathrm\{mph\}\textbackslash{}). The pilot forgot to check the weather and doesn't know that the wind is blowing to the south at \textbackslash{}(50\textbackslash{},\textbackslash{}mathrm\{mph\}\textbackslash{})

\#\# Question

What is the plane's ground speed?

\#\# Answer choices

- A: A: 510pmh

- B: B: 502mph

- C: C: 490pmh

- D: D: 498pmh

\#\# Figure caption (auxiliary; use the image as primary evidence)

The image is a vector diagram illustrating motion with respect to air and ground. It involves two main locations: Chicago and Cleveland, indicated by black dots.

There are three vectors:
1. **\textbackslash{}(\textbackslash{}vec\{v\}\_\{PA\}\textbackslash{})**: Represents the velocity of the plane relative to the air. It is a horizontal arrow from Chicago towards a plane above the line connecting the two cities.
2. **\textbackslash{}(\textbackslash{}vec\{v\}\_\{PG\}\textbackslash{})**: Represents the velocity of the plane relative to the ground. It connects Chicago to a plane on a diagonal path toward Cleveland.
3. **\textbackslash{}(\textbackslash{}vec\{v\}\_\{AG\}\textbackslash{})**: Represents the velocity of the air. It forms a vertical line connected from the endpoint of \textbackslash{}(\textbackslash{}vec\{v\}\_\{PA\}\textbackslash{}) to the endpoint of \textbackslash{}(\textbackslash{}vec\{v\}\_\{PG\}\textbackslash{}).

Two illustrations of planes indicate directions along \textbackslash{}(\textbackslash{}vec\{v\}\_\{PA\}\textbackslash{}) and \textbackslash{}(\textbackslash{}vec\{v\}\_\{PG\}\textbackslash{}).

Text labels are present for each vector to describe what they represent, such as "of plane relative to air" or "of air."

\#\# Dataset metadata

Kinematics; Spatial Relation Reasoning, Multi-Formula Reasoning

\paragraph{Recorded answer/context.}
B: B: 502mph

\paragraph{Figure/image path.}
/root/proposal\_for\_physic/science-mango/phyx\_mini\_run/phyx\_data/test\_image/828.png

\paragraph{Formalization target.}
create a compiling Lean file with sorry bodies at `PhyXMiniProblems/problem\_phyx\_mini\_0828.lean`. The Lean declarations must preserve the physical quantities, dimensions or dimensional roles, figure labels, governing-law hypotheses, and final relation expressed by this problem.
Use Mathlib/Physlib names found through LeanExplore where available. If a physics API is missing, introduce faithful local abstractions rather than scalar placeholder aliases.

\begin{theorem}[Physics formalization target]
\label{thm:physics:phyx_mini_0828:target}
\lean{PhyXMiniProblems.ProblemPhyXMini0828.problem_phyx_mini_0828}
\uses{def:physics:phyx-mini-0828:phyxminiproblems-problemphyxmini0828-speedinmilesperhour, def:physics:phyx-mini-0828:phyxminiproblems-problemphyxmini0828-planewindsetup, def:physics:phyx-mini-0828:phyxminiproblems-problemphyxmini0828-matchesscenarioandsuppliedfigure, def:physics:phyx-mini-0828:phyxminiproblems-problemphyxmini0828-statedplanewinddata, def:physics:phyx-mini-0828:phyxminiproblems-problemphyxmini0828-planewindkinematicslaws, def:physics:phyx-mini-0828:phyxminiproblems-problemphyxmini0828-answerchoicemilesperhour}
The assigned autoformalize agent should translate this physics problem into Lean declarations in the covered file, with theorem and lemma proof bodies written as `by sorry`.
\end{theorem}
\begin{proof}
This is an autoformalization task, not a proof task. Produce statements that can later be proved by the physics prover without weakening the physical model.
\end{proof}

% --- Archon declaration topology begin ---
\section{Declaration topology}
\begin{definition}[velocity Dimension]
  \label{def:physics:phyx-mini-0828:phyxminiproblems-problemphyxmini0828-velocitydimension}
  \lean{PhyXMiniProblems.ProblemPhyXMini0828.velocityDimension}
  Velocity has physical dimension length divided by time
\end{definition}
\begin{proof}
  This is the declared model definition of velocity Dimension.
\end{proof}
\begin{definition}[Plane Readout]
  \label{def:physics:phyx-mini-0828:phyxminiproblems-problemphyxmini0828-planereadout}
  \lean{PhyXMiniProblems.ProblemPhyXMini0828.PlaneReadout}
  The Euclidean plane in which the city positions and velocity components are read
\end{definition}
\begin{proof}
  This is the declared model definition of Plane Readout.
\end{proof}
\begin{definition}[Length Quantity]
  \label{def:physics:phyx-mini-0828:phyxminiproblems-problemphyxmini0828-lengthquantity}
  \lean{PhyXMiniProblems.ProblemPhyXMini0828.LengthQuantity}
  A nonnegative, unit-independent physical distance
\end{definition}
\begin{proof}
  This is the declared model definition of Length Quantity.
\end{proof}
\begin{definition}[Speed Quantity]
  \label{def:physics:phyx-mini-0828:phyxminiproblems-problemphyxmini0828-speedquantity}
  \lean{PhyXMiniProblems.ProblemPhyXMini0828.SpeedQuantity}
  A nonnegative, unit-independent speed magnitude supplied by Physlib
\end{definition}
\begin{proof}
  This is the declared model definition of Speed Quantity.
\end{proof}
\begin{definition}[Planar Position Quantity]
  \label{def:physics:phyx-mini-0828:phyxminiproblems-problemphyxmini0828-planarpositionquantity}
  \lean{PhyXMiniProblems.ProblemPhyXMini0828.PlanarPositionQuantity}
  \uses{def:physics:phyx-mini-0828:phyxminiproblems-problemphyxmini0828-planereadout}
  A unit-independent signed planar position
\end{definition}
\begin{proof}
  This is the declared model definition of Planar Position Quantity.
\end{proof}
\begin{definition}[Planar Velocity Quantity]
  \label{def:physics:phyx-mini-0828:phyxminiproblems-problemphyxmini0828-planarvelocityquantity}
  \lean{PhyXMiniProblems.ProblemPhyXMini0828.PlanarVelocityQuantity}
  \uses{def:physics:phyx-mini-0828:phyxminiproblems-problemphyxmini0828-velocitydimension, def:physics:phyx-mini-0828:phyxminiproblems-problemphyxmini0828-planereadout}
  A unit-independent signed planar velocity
\end{definition}
\begin{proof}
  This is the declared model definition of Planar Velocity Quantity.
\end{proof}
\begin{definition}[length Readout]
  \label{def:physics:phyx-mini-0828:phyxminiproblems-problemphyxmini0828-lengthreadout}
  \lean{PhyXMiniProblems.ProblemPhyXMini0828.lengthReadout}
  \uses{def:physics:phyx-mini-0828:phyxminiproblems-problemphyxmini0828-lengthquantity}
  Read a physical distance in a selected length unit
\end{definition}
\begin{proof}
  This is the declared model definition of length Readout.
\end{proof}
\begin{definition}[speed Readout]
  \label{def:physics:phyx-mini-0828:phyxminiproblems-problemphyxmini0828-speedreadout}
  \lean{PhyXMiniProblems.ProblemPhyXMini0828.speedReadout}
  \uses{def:physics:phyx-mini-0828:phyxminiproblems-problemphyxmini0828-speedquantity}
  Read a speed in coherent selected length and time units
\end{definition}
\begin{proof}
  This is the declared model definition of speed Readout.
\end{proof}
\begin{definition}[planar Position Readout]
  \label{def:physics:phyx-mini-0828:phyxminiproblems-problemphyxmini0828-planarpositionreadout}
  \lean{PhyXMiniProblems.ProblemPhyXMini0828.planarPositionReadout}
  \uses{def:physics:phyx-mini-0828:phyxminiproblems-problemphyxmini0828-planereadout, def:physics:phyx-mini-0828:phyxminiproblems-problemphyxmini0828-planarpositionquantity}
  Read a planar position in a selected length unit
\end{definition}
\begin{proof}
  This is the declared model definition of planar Position Readout.
\end{proof}
\begin{definition}[planar Velocity Readout]
  \label{def:physics:phyx-mini-0828:phyxminiproblems-problemphyxmini0828-planarvelocityreadout}
  \lean{PhyXMiniProblems.ProblemPhyXMini0828.planarVelocityReadout}
  \uses{def:physics:phyx-mini-0828:phyxminiproblems-problemphyxmini0828-planereadout, def:physics:phyx-mini-0828:phyxminiproblems-problemphyxmini0828-planarvelocityquantity}
  Read a planar velocity in coherent selected length and time units
\end{definition}
\begin{proof}
  This is the declared model definition of planar Velocity Readout.
\end{proof}
\begin{definition}[length In Miles]
  \label{def:physics:phyx-mini-0828:phyxminiproblems-problemphyxmini0828-lengthinmiles}
  \lean{PhyXMiniProblems.ProblemPhyXMini0828.lengthInMiles}
  \uses{def:physics:phyx-mini-0828:phyxminiproblems-problemphyxmini0828-lengthquantity, def:physics:phyx-mini-0828:phyxminiproblems-problemphyxmini0828-lengthreadout}
  Mile readout used for the Chicago--Cleveland separation
\end{definition}
\begin{proof}
  This is the declared model definition of length In Miles.
\end{proof}
\begin{definition}[speed In Miles Per Hour]
  \label{def:physics:phyx-mini-0828:phyxminiproblems-problemphyxmini0828-speedinmilesperhour}
  \lean{PhyXMiniProblems.ProblemPhyXMini0828.speedInMilesPerHour}
  \uses{def:physics:phyx-mini-0828:phyxminiproblems-problemphyxmini0828-speedquantity, def:physics:phyx-mini-0828:phyxminiproblems-problemphyxmini0828-speedreadout}
  Mile-per-hour readout used for all speeds in the problem and choices
\end{definition}
\begin{proof}
  This is the declared model definition of speed In Miles Per Hour.
\end{proof}
\begin{definition}[position In Miles]
  \label{def:physics:phyx-mini-0828:phyxminiproblems-problemphyxmini0828-positioninmiles}
  \lean{PhyXMiniProblems.ProblemPhyXMini0828.positionInMiles}
  \uses{def:physics:phyx-mini-0828:phyxminiproblems-problemphyxmini0828-planereadout, def:physics:phyx-mini-0828:phyxminiproblems-problemphyxmini0828-planarpositionquantity, def:physics:phyx-mini-0828:phyxminiproblems-problemphyxmini0828-planarpositionreadout}
  Planar position read in miles, with east as positive x and north as positive y
\end{definition}
\begin{proof}
  This is the declared model definition of position In Miles.
\end{proof}
\begin{definition}[velocity In Miles Per Hour]
  \label{def:physics:phyx-mini-0828:phyxminiproblems-problemphyxmini0828-velocityinmilesperhour}
  \lean{PhyXMiniProblems.ProblemPhyXMini0828.velocityInMilesPerHour}
  \uses{def:physics:phyx-mini-0828:phyxminiproblems-problemphyxmini0828-planereadout, def:physics:phyx-mini-0828:phyxminiproblems-problemphyxmini0828-planarvelocityquantity, def:physics:phyx-mini-0828:phyxminiproblems-problemphyxmini0828-planarvelocityreadout}
  Planar velocity read in miles per hour in the same east/north coordinates
\end{definition}
\begin{proof}
  This is the declared model definition of velocity In Miles Per Hour.
\end{proof}
\begin{definition}[City Label]
  \label{def:physics:phyx-mini-0828:phyxminiproblems-problemphyxmini0828-citylabel}
  \lean{PhyXMiniProblems.ProblemPhyXMini0828.CityLabel}
  The two cities named in the written scenario and drawn as black dots
\end{definition}
\begin{proof}
  This is the declared model definition of City Label.
\end{proof}
\begin{definition}[Physical Body]
  \label{def:physics:phyx-mini-0828:phyxminiproblems-problemphyxmini0828-physicalbody}
  \lean{PhyXMiniProblems.ProblemPhyXMini0828.PhysicalBody}
  Bodies used to interpret the two-letter subscripts on relative velocities
\end{definition}
\begin{proof}
  This is the declared model definition of Physical Body.
\end{proof}
\begin{definition}[Relative Velocity Label]
  \label{def:physics:phyx-mini-0828:phyxminiproblems-problemphyxmini0828-relativevelocitylabel}
  \lean{PhyXMiniProblems.ProblemPhyXMini0828.RelativeVelocityLabel}
  The three literal vector labels printed in primary image 828.png
\end{definition}
\begin{proof}
  This is the declared model definition of Relative Velocity Label.
\end{proof}
\begin{definition}[moving Body]
  \label{def:physics:phyx-mini-0828:phyxminiproblems-problemphyxmini0828-movingbody}
  \lean{PhyXMiniProblems.ProblemPhyXMini0828.movingBody}
  \uses{def:physics:phyx-mini-0828:phyxminiproblems-problemphyxmini0828-physicalbody, def:physics:phyx-mini-0828:phyxminiproblems-problemphyxmini0828-relativevelocitylabel}
  The moving body denoted by the first letter of each velocity subscript
\end{definition}
\begin{proof}
  This is the declared model definition of moving Body.
\end{proof}
\begin{definition}[reference Body]
  \label{def:physics:phyx-mini-0828:phyxminiproblems-problemphyxmini0828-referencebody}
  \lean{PhyXMiniProblems.ProblemPhyXMini0828.referenceBody}
  \uses{def:physics:phyx-mini-0828:phyxminiproblems-problemphyxmini0828-physicalbody, def:physics:phyx-mini-0828:phyxminiproblems-problemphyxmini0828-relativevelocitylabel}
  The reference body denoted by the second letter of each velocity subscript
\end{definition}
\begin{proof}
  This is the declared model definition of reference Body.
\end{proof}
\begin{definition}[Cardinal Direction]
  \label{def:physics:phyx-mini-0828:phyxminiproblems-problemphyxmini0828-cardinaldirection}
  \lean{PhyXMiniProblems.ProblemPhyXMini0828.CardinalDirection}
  Cardinal directions needed to transcribe the prose data
\end{definition}
\begin{proof}
  This is the declared model definition of Cardinal Direction.
\end{proof}
\begin{definition}[Pilot Wind Knowledge]
  \label{def:physics:phyx-mini-0828:phyxminiproblems-problemphyxmini0828-pilotwindknowledge}
  \lean{PhyXMiniProblems.ProblemPhyXMini0828.PilotWindKnowledge}
  The pilot's stated knowledge of the wind before departure
\end{definition}
\begin{proof}
  This is the declared model definition of Pilot Wind Knowledge.
\end{proof}
\begin{definition}[Figure Point]
  \label{def:physics:phyx-mini-0828:phyxminiproblems-problemphyxmini0828-figurepoint}
  \lean{PhyXMiniProblems.ProblemPhyXMini0828.FigurePoint}
  Distinguished endpoints in the head-to-tail vector diagram
\end{definition}
\begin{proof}
  This is the declared model definition of Figure Point.
\end{proof}
\begin{definition}[Figure Object]
  \label{def:physics:phyx-mini-0828:phyxminiproblems-problemphyxmini0828-figureobject}
  \lean{PhyXMiniProblems.ProblemPhyXMini0828.FigureObject}
  Physical and graphical objects visible in the supplied image
\end{definition}
\begin{proof}
  This is the declared model definition of Figure Object.
\end{proof}
\begin{definition}[Figure Text Label]
  \label{def:physics:phyx-mini-0828:phyxminiproblems-problemphyxmini0828-figuretextlabel}
  \lean{PhyXMiniProblems.ProblemPhyXMini0828.FigureTextLabel}
  Literal city and vector-description text printed in the supplied image
\end{definition}
\begin{proof}
  This is the declared model definition of Figure Text Label.
\end{proof}
\begin{definition}[Supplied Velocity Diagram]
  \label{def:physics:phyx-mini-0828:phyxminiproblems-problemphyxmini0828-suppliedvelocitydiagram}
  \lean{PhyXMiniProblems.ProblemPhyXMini0828.SuppliedVelocityDiagram}
  \uses{def:physics:phyx-mini-0828:phyxminiproblems-problemphyxmini0828-planarvelocityquantity, def:physics:phyx-mini-0828:phyxminiproblems-problemphyxmini0828-relativevelocitylabel, def:physics:phyx-mini-0828:phyxminiproblems-problemphyxmini0828-figurepoint, def:physics:phyx-mini-0828:phyxminiproblems-problemphyxmini0828-figureobject, def:physics:phyx-mini-0828:phyxminiproblems-problemphyxmini0828-figuretextlabel}
  Typed primary-image data. The diagram's three vector values remain physical planar velocities. Arrow endpoints record the head-to-tail geometry without postulating any numerical value for vPG
\end{definition}
\begin{proof}
  This is the declared model definition of Supplied Velocity Diagram.
\end{proof}
\begin{definition}[Plane Wind Setup]
  \label{def:physics:phyx-mini-0828:phyxminiproblems-problemphyxmini0828-planewindsetup}
  \lean{PhyXMiniProblems.ProblemPhyXMini0828.PlaneWindSetup}
  \uses{def:physics:phyx-mini-0828:phyxminiproblems-problemphyxmini0828-lengthquantity, def:physics:phyx-mini-0828:phyxminiproblems-problemphyxmini0828-speedquantity, def:physics:phyx-mini-0828:phyxminiproblems-problemphyxmini0828-planarpositionquantity, def:physics:phyx-mini-0828:phyxminiproblems-problemphyxmini0828-planarvelocityquantity, def:physics:phyx-mini-0828:phyxminiproblems-problemphyxmini0828-citylabel, def:physics:phyx-mini-0828:phyxminiproblems-problemphyxmini0828-cardinaldirection, def:physics:phyx-mini-0828:phyxminiproblems-problemphyxmini0828-pilotwindknowledge, def:physics:phyx-mini-0828:phyxminiproblems-problemphyxmini0828-suppliedvelocitydiagram}
  Independent physical observables in the written problem. In particular, planeVelocityRelativeToGround and planeGroundSpeed are unconstrained here; their relation to the given velocities is supplied only by governing laws
\end{definition}
\begin{proof}
  This is the declared model definition of Plane Wind Setup.
\end{proof}
\begin{definition}[Matches Scenario And Supplied Figure]
  \label{def:physics:phyx-mini-0828:phyxminiproblems-problemphyxmini0828-matchesscenarioandsuppliedfigure}
  \lean{PhyXMiniProblems.ProblemPhyXMini0828.MatchesScenarioAndSuppliedFigure}
  \uses{def:physics:phyx-mini-0828:phyxminiproblems-problemphyxmini0828-planewindsetup}
  Qualitative prose and primary-image facts. These fields identify the three pictured velocity vectors but do not give the magnitude of vPG or the plane's ground speed
\end{definition}
\begin{proof}
  This is the declared model definition of Matches Scenario And Supplied Figure.
\end{proof}
\begin{definition}[Stated Plane Wind Data]
  \label{def:physics:phyx-mini-0828:phyxminiproblems-problemphyxmini0828-statedplanewinddata}
  \lean{PhyXMiniProblems.ProblemPhyXMini0828.StatedPlaneWindData}
  \uses{def:physics:phyx-mini-0828:phyxminiproblems-problemphyxmini0828-lengthinmiles, def:physics:phyx-mini-0828:phyxminiproblems-problemphyxmini0828-speedinmilesperhour, def:physics:phyx-mini-0828:phyxminiproblems-problemphyxmini0828-positioninmiles, def:physics:phyx-mini-0828:phyxminiproblems-problemphyxmini0828-velocityinmilesperhour, def:physics:phyx-mini-0828:phyxminiproblems-problemphyxmini0828-planewindsetup}
  The numerical values and directions explicitly given in the prose. East is positive x and north is positive y, so a southward component is negative. No field assigns a component or magnitude to the ground-relative velocity
\end{definition}
\begin{proof}
  This is the declared model definition of Stated Plane Wind Data.
\end{proof}
\begin{definition}[Plane Wind Kinematics Laws]
  \label{def:physics:phyx-mini-0828:phyxminiproblems-problemphyxmini0828-planewindkinematicslaws}
  \lean{PhyXMiniProblems.ProblemPhyXMini0828.PlaneWindKinematicsLaws}
  \uses{def:physics:phyx-mini-0828:phyxminiproblems-problemphyxmini0828-lengthreadout, def:physics:phyx-mini-0828:phyxminiproblems-problemphyxmini0828-speedreadout, def:physics:phyx-mini-0828:phyxminiproblems-problemphyxmini0828-planarpositionreadout, def:physics:phyx-mini-0828:phyxminiproblems-problemphyxmini0828-planarvelocityreadout, def:physics:phyx-mini-0828:phyxminiproblems-problemphyxmini0828-planewindsetup}
  Governing geometry and Galilean kinematics. Each law is stated for every coherent choice of length and time units. The final numerical ground speed is not a law field: it must be derived by adding vPA and vAG and taking the Euclidean norm of vPG
\end{definition}
\begin{proof}
  This is the declared model definition of Plane Wind Kinematics Laws.
\end{proof}
\begin{definition}[Answer Choice]
  \label{def:physics:phyx-mini-0828:phyxminiproblems-problemphyxmini0828-answerchoice}
  \lean{PhyXMiniProblems.ProblemPhyXMini0828.AnswerChoice}
  The four answer-choice labels from the problem statement
\end{definition}
\begin{proof}
  This is the declared model definition of Answer Choice.
\end{proof}
\begin{definition}[answer Choice Miles Per Hour]
  \label{def:physics:phyx-mini-0828:phyxminiproblems-problemphyxmini0828-answerchoicemilesperhour}
  \lean{PhyXMiniProblems.ProblemPhyXMini0828.answerChoiceMilesPerHour}
  \uses{def:physics:phyx-mini-0828:phyxminiproblems-problemphyxmini0828-answerchoice}
  The numerical mile-per-hour values printed beside the choices. The source's pmh spellings in choices A and C are interpreted as the evident mph typo
\end{definition}
\begin{proof}
  This is the declared model definition of answer Choice Miles Per Hour.
\end{proof}
% --- Archon declaration topology end ---
% --- Archon physics formalization source end ---
